% Source: https://artofproblemsolving.com/community/c3h2574373_odd_numbers
% Stable identity: aops-topic-2574373-q01
If $n=2m$ is even, the odd numbers are $1,3,\dots,2m-1$, so there are $m=\dfrac n2$ of them. If $n=2m+1$ is odd, they are $1,3,\dots,2m+1$, so there are $m+1=\dfrac{n+1}2$. In both cases the count is
\[\boxed{\left\lceil\dfrac n2\right\rceil=\left\lfloor\dfrac{n+1}2\right\rfloor}.\]
The claim $\left\lfloor\dfrac n2\right\rfloor$ is correct only for even $n$; for odd $n$ it is one too small (for $n=1$ it gives $0$, but $1$ is odd). In fact $\left\lfloor\dfrac n2\right\rfloor$ counts the \emph{even} numbers from $1$ to $n$.
